% Paper 2, Section 5 — Copy: the precise-recall workload (evidence E3).
% Numbers regenerated from `python agg_copy.py 30` at drafting time and verified against ledger
% entries a8e7a7c (M=16), 8e74f48 (M=32), 159e341 (M=32 floor control), 5bb7bcc/8d050cb/48510c8
% (regularization control + corrected margins). M=64 column from the completed ep30 row.
% CLAIMS DISCIPLINE: M=16/32 margins are quoted against the NOISE-REGULARIZED reference; the
% M=64 column has NO regularized reference (the control ran at L in {33,65} only) and the
% digital baseline is near-floor there — M=64 is reported for completeness, not used for claims.

\section{Copy: the precise-recall workload}
\label{sec:copy}

The copy task inverts every conclusion of Section~\ref{sec:charlm}, and the inversion is the
pivot of this paper. The model must reproduce $M \in \{16, 32, 64\}$ symbols from a $K{=}16$
alphabet after a delay (chance accuracy $0.0625$, chance bpc $4.000$); unlike character-level
language modeling, success requires \emph{precise retention} of specific symbols rather than
sequence statistics.

\subsection{Results and the ranking inversion}
\label{sec:copy-results}

Table~\ref{tab:copy} gives the three-seed comparison at all three memory loads. ``Margin
kept'' is the fraction of the reference baseline's above-chance accuracy margin a variant
retains --- the appropriate metric when absolute accuracies are low. Following
Section~\ref{sec:controls}, the reference at $M{=}16$ and $M{=}32$ is the
\emph{noise-regularized} digital baseline, which is stronger than the plain one on copy as
well (paired $\Delta$bpc $-0.150 \pm 0.033$ at $M{=}16$, $-0.215 \pm 0.011$ at $M{=}32$);
margins against the plain baseline are $3$--$5$ points higher and are upper bounds.

\begin{table*}[t]
\centering
\caption{Copy task, 3 seeds (mean $\pm$ sd), ep30 budget, $H{=}256$ ($87{,}889$ parameters),
$K{=}16$, chance acc $0.0625$ / bpc $4.000$. ``Emit'' is the emitted (wire) rate; ``state''
is state activity. Margin kept is vs.\ the noise-regularized digital reference at
$M{=}16/32$ and vs.\ the plain digital baseline at $M{=}64$ (the regularization control was
not run there; see \S\ref{sec:copy-load}). Energy is the 45\,nm proxy, not a silicon
measurement.}
\label{tab:copy}
\footnotesize
\begin{tabular}{l l r r r r r r}
\toprule
$M$ & Variant & acc & bpc & emit & state & margin kept & pJ/tok \\
\midrule
16 & digital (plain)      & $0.630 \pm 0.011$ & $1.420 \pm 0.034$ & 1.000 & 1.000 & 0.95 & 398k \\
16 & digital + noise reg. & $0.663 \pm 0.009$ & $1.270 \pm 0.029$ & 1.000 & 1.000 & 1.00 & 398k \\
16 & spikeout             & $0.583 \pm 0.012$ & $1.545 \pm 0.071$ & 0.447 & 1.000 & 0.87 & 123k \\
16 & analog $\theta{=}0.2$ & $0.361 \pm 0.051$ & $2.777 \pm 0.226$ & 0.501 & 0.972 & 0.50 & 247k \\
16 & spikestate           & $0.107 \pm 0.016$ & $3.929 \pm 0.037$ & 0.496 & 0.496 & 0.074 & 108k \\
\midrule
32 & digital (plain)      & $0.315 \pm 0.003$ & $2.893 \pm 0.013$ & 1.000 & 1.000 & 0.90 & 398k \\
32 & digital + noise reg. & $0.343 \pm 0.003$ & $2.678 \pm 0.016$ & 1.000 & 1.000 & 1.00 & 398k \\
32 & spikeout             & $0.180 \pm 0.001$ & $3.605 \pm 0.003$ & 0.485 & 1.000 & 0.42 & 125k \\
32 & analog $\theta{=}0.2$ & $0.142 \pm 0.000$ & $3.788 \pm 0.002$ & 0.449 & 0.964 & 0.28 & 231k \\
32 & spikestate           & $0.062 \pm 0.001$ & $4.005 \pm 0.004$ & 0.405 & 0.405 & 0.00 & 102k \\
\midrule
64 & digital (plain)      & $0.167 \pm 0.001$ & $3.666 \pm 0.012$ & 1.000 & 1.000 & 1.00 & 398k \\
64 & spikeout             & $0.133 \pm 0.001$ & $3.824 \pm 0.002$ & 0.481 & 1.000 & 0.67 & 125k \\
64 & analog $\theta{=}0.2$ & $0.110 \pm 0.001$ & $3.914 \pm 0.001$ & 0.401 & 0.955 & 0.46 & 216k \\
64 & spikestate           & $0.062 \pm 0.001$ & $4.002 \pm 0.001$ & 0.238 & 0.238 & 0.00 & 91k \\
\bottomrule
\end{tabular}
\end{table*}

The three non-digital variants emit at nearly matched rates ($0.40$--$0.50$ at every load), so
--- as in Section~\ref{sec:charlm-matched} --- the quality spread is not an operating-point
artifact. And the ordering is the exact reverse of char-LM: at $M{=}16$, best to worst is
\textbf{spikeout} (margin kept $0.87$), \textbf{analog} ($0.50$), \textbf{spikestate}
($0.074$); on char-LM the same three variants ranked analog ($\Delta$bpc $+0.160$ vs.\ the
regularized reference) $\ll$ spikestate ($+0.587$) $\ll$ spikeout ($+1.364$). The variant
that was worst on the statistical workload --- continuous state, spiked output --- is the
only one that survives precise recall, keeping $87\%$ of the regularized baseline's margin at
a $45\%$ emission rate and $3.2\times$ lower proxy energy.

\subsection{A one-bit recurrent state eliminates recall}
\label{sec:copy-floor}

The spikestate row is not merely degraded. At $M{=}32$ it lands at \emph{exactly} chance ---
accuracy $0.0623 \pm 0.0012$ against $0.0625$, bpc $4.005 \pm 0.004$ against $4.000$ --- and
reproduces this at $M{=}64$ ($0.0623 \pm 0.0005$). A one-bit recurrent state does not degrade
precise-recall performance at these loads; it \emph{eliminates} it. It is simultaneously the
cheapest cell in the table ($102$k pJ/token at $M{=}32$, $4\times$ below the digital
baseline) and buys nothing: the cleanest demonstration in this study that an energy number is
meaningless without the task quality it purchased. For memory-bearing workloads, carrying the
recurrent state in spikes is a non-option, not a cheap option.

\subsection{Load dependence, and the validity limit at $M{=}64$}
\label{sec:copy-load}

Between $M{=}16$ and $M{=}32$ --- the loads where the digital reference still learns the task
well above chance --- the quality gap \emph{widens} for both surviving neuromorphic routes:
margin kept falls $0.87 \rightarrow 0.42$ for spikeout and $0.50 \rightarrow 0.28$ for analog
(seed sd $\le 0.003$ acc, so far outside noise). The analog variant's emission does not rise
to compensate at fixed $\theta$; it \emph{falls} ($0.501 \rightarrow 0.449$). Two
consequences: there is no single ``$X\%$ of digital'' figure for a neuromorphic SSM --- the
workload's retention demand sets it --- and any analog energy/quality curve must be drawn per
memory load.

The $M{=}64$ column is reported for completeness but is \emph{validity-limited}, and we do
not use it for claims. The digital baseline itself is close to floor there (accuracy $0.167$,
only $2.7\times$ chance, against $10\times$ chance at $M{=}16$), so every variant compresses
toward chance and both metrics lose resolution: margin kept rebounds non-monotonically
(spikeout $0.92 \rightarrow 0.46 \rightarrow 0.67$ vs.\ the plain baseline) and paired
$\Delta$bpc compresses ($+0.71 \rightarrow +0.16$ for spikeout from $M{=}32$ to $M{=}64$) ---
a shrinking-reference artifact, not a recovery. The regularization control was also not run
at $M{=}64$, so no corrected margin exists there. The load-dependence claim above is
therefore made on $M{=}16 \rightarrow 32$ only.
